
% ============================================================
% APÉNDICE A3 — Figuras con datos REALES (Gasolina Super)
% Generado automáticamente desde los outputs del modelo TCROC-Markov
% ============================================================

\section{Ilustraciones del Proceso Estocástico \texorpdfstring{$\{S_t\}$}{St}}
\label{apx:figuras_proceso_st}

\begin{figure}[htbp]
\centering
\begin{tikzpicture}
  \begin{axis}[
    width=0.8\textwidth,
    xlabel={Tiempo \( t \)}, ylabel={Estado \( S_t \)},
    ymin=0.5, ymax=4.5,
    ytick={1,2,3,4}, 
    yticklabels={\\textnormal{s\\_1},\\textnormal{s\\_2},\\textnormal{s\\_3},\\textnormal{s\\_4}},
    legend pos=north west,
    grid=major,
    title={Trayectorias del proceso $\{S_t\}$ (Gasolina Super)}
  ]
  \addplot+[mark=*, thick, nordFrost] coordinates {
    (1,2) (2,3) (3,3) (4,3) (5,3) (6,3) (7,3) (8,3) (9,3) (10,3)
  };
  \addplot+[mark=triangle*, thick, nordRed] coordinates {
    (1,2) (2,2) (3,2) (4,2) (5,2) (6,2) (7,2) (8,2) (9,2) (10,2)
  };
  \legend{Observada, Simulada ($\hat{P}$)}
  \end{axis}
\end{tikzpicture}
\caption{Trayectorias del proceso $\{S_t\}$ para la Gasolina~Super: 
la secuencia observada (datos reales del período 2017~--~2025) frente a 
una realización simulada a partir de la matriz $\hat{P}$ estimada 
por NNLS.}
\label{fig:apx_st_trayectorias}
\end{figure}

\begin{figure}[htbp]
\centering
\begin{tikzpicture}[node distance=2cm and 2.5cm, >=Stealth]
\node[state, fill=nordRed!15] (s1)  {\( s1 \)};
\node[state, fill=nordGreen!15] (s2) [right=of s1] {\( s2 \)};
\node[state, fill=nordYellow!15] (s3) [below=of s2] {\( s3 \)};
\node[state, fill=nordFrost!15] (s4) [left=of s3] {\( s4 \)};

\draw[->] (s1) to[loop above] node {0.90} (s1);
\draw[->] (s1) to[bend left=15] node[right] {0.10} (s2);
\draw[->] (s2) to[loop above] node {0.92} (s2);
\draw[->] (s3) to[bend left=15] node[above] {0.08} (s2);
\draw[->] (s3) to[loop below] node {0.89} (s3);
\draw[->] (s4) to[loop below] node {0.95} (s4);
\end{tikzpicture}
\caption{Diagrama de transiciones entre estados $s_j$ para la 
Gasolina~Super con $K=4$ regímenes. Las probabilidades 
corresponden a la matriz $\hat{P}$ estimada por NNLS 
(véase Capítulo~\ref{cap:modelo_hibrido}).}
\label{fig:apx_diagrama_transiciones}
\end{figure}

\vspace{1em}
\noindent\textbf{Nota:} Las figuras de esta sección están generadas 
a partir de las matrices $\hat{P}$ estimadas con datos reales de 
precios de combustibles de Honduras (2017--2025). La trayectoria 
observada corresponde a la serie de la Gasolina~Super procesada con 
los hiperparámetros óptimos obtenidos en el Capítulo~\ref{cap:regimenes_combustibles}.
